\section{Resolution robustness of the geopotential dispersion}
\label{sec:varrobust}

The difficulty identified in Section~\ref{sec:slopesens} is not that the slope
is noisy---it is smooth and monotonic in the mesh---but that it does not
approach a limit. This section shows that the normalized geopotential dispersion
$\potvarn$ of Eq.~\eqref{eq:potvarn}, computed from the same fields on the same
meshes, behaves differently: it too rises with refinement, but it saturates.

\subsection{The measurement}
\label{sec:varmeas}

The measurement is identical to that of Section~\ref{sec:slopemeas} and is
performed on the same decimation sequences, from the same evaluations. No
additional solver calls are involved: $\slopeavg$ and $\potvarn$ are two
summaries of one computed field, so any difference between them is a difference
between the diagnostics, not between the calculations that produced them.

\subsection{Results}
\label{sec:varres}

Table~\ref{tab:varseries} gives the sequences, and Figure~\ref{fig:resolution}
plots them directly beneath the corresponding slope curves, on identical axes.
As with the slope, $\potvarn$ increases monotonically with facet count on every
body. Unlike the slope, it
flattens: on Itokawa and Gaspra the last two meshes return the same value to
five significant figures, and on Ida they differ in the fifth.

\begin{table}[t]
\centering
\small
\caption{Normalized geopotential dispersion $\potvarn$ as a function of facet
count, on the same decimation sequences as Table~\ref{tab:slopeseries}.}
\label{tab:varseries}
\begin{tabular}{lrr}
\toprule
Body & $\Nf$ & $\potvarn$ \\
\midrule
Ryugu   & $3\,096$  & $0.03153$ \\
        & $5\,780$  & $0.03233$ \\
        & $11\,134$ & $0.03295$ \\
        & $20\,014$ & $0.03327$ \\
        & $42\,834$ & $0.03347$ \\
\midrule
Eros    & $2\,097$  & $0.03241$ \\
        & $6\,306$  & $0.03345$ \\
        & $12\,155$ & $0.03384$ \\
        & $21\,850$ & $0.03403$ \\
        & $36\,030$ & $0.03413$ \\
\midrule
Itokawa & $2\,183$  & $0.08527$ \\
        & $6\,404$  & $0.08577$ \\
        & $12\,272$ & $0.08598$ \\
        & $21\,196$ & $0.08608$ \\
        & $41\,369$ & $0.08608$ \\
\midrule
Gaspra  & $2\,821$  & $0.05929$ \\
        & $5\,240$  & $0.05980$ \\
        & $12\,500$ & $0.06009$ \\
        & $20\,773$ & $0.06017$ \\
        & $32\,040$ & $0.06017$ \\
\midrule
Ida     & $2\,141$  & $0.03490$ \\
        & $5\,808$  & $0.03572$ \\
        & $11\,804$ & $0.03626$ \\
        & $20\,673$ & $0.03641$ \\
        & $32\,040$ & $0.03643$ \\
\bottomrule
\end{tabular}
\end{table}

\begin{figure}[tp]
\centering
\includegraphics[width=\textwidth]{fig3.pdf}
\caption{Behaviour of the two diagnostics under mesh refinement, expressed as
the fractional deviation from the value obtained on the finest mesh of each
sequence. (a) Area-weighted mean slope. (b) Normalized geopotential dispersion,
plotted on \emph{identical} axes. Both quantities are computed from the same
field on the same meshes, so the contrast is a property of the diagnostics
rather than of the calculation. The slope curves spread by up to a quarter of
their own value and are still moving at the finest meshes; the dispersion curves
are flat on this scale.}
\label{fig:resolution}
\end{figure}

\subsection{The two diagnostics compared}
\label{sec:varcompare}

Table~\ref{tab:varsens} and Figure~\ref{fig:diagnostics} place the two
diagnostics side by side. We report both
measures of Section~\ref{sec:diagnostics}, because they support claims of
different strength and it would be misleading to quote only the more favourable
one.

By the range sensitivity $S$, the geopotential dispersion is the more stable
diagnostic on every body, but the margin varies considerably: it is a factor of
$13.6$ on Itokawa, $4.3$ on Ryugu and $3.6$ on Eros, but only $1.3$ on Ida and
$1.2$ on Gaspra. The last two figures should not be read as a weakness of the
dispersion. Gaspra and Ida are precisely the bodies whose slope was already
stable ($S[\slopeavg]=1.8\%$ and $5.6\%$), so there was little for a better
diagnostic to improve upon; the ratio is uninformative where the denominator is
already small. The honest summary of the range metric is that $\potvarn$ varies
by $0.9\%$ to $5.9\%$ where $\slopeavg$ varies by $1.8\%$ to $25.6\%$, and that
the advantage is largest exactly where it matters most.

\begin{table}[t]
\centering
\small
\caption{Range sensitivity $S$ and convergence residual $C$ for the two
diagnostics, on identical meshes. A dash indicates that the diagnostic did not
change over the final refinement step at the precision retained.}
\label{tab:varsens}
\begin{tabular}{lrrrr}
\toprule
& \multicolumn{2}{c}{$S$} & \multicolumn{2}{c}{$C$} \\
\cmidrule(lr){2-3}\cmidrule(lr){4-5}
Body & $\slopeavg$ & $\potvarn$ & $\slopeavg$ & $\potvarn$ \\
\midrule
Ryugu   & $25.6\%$ & $5.9\%$ & $8.24\%$ & $0.60\%$ \\
Eros    & $18.3\%$ & $5.1\%$ & $4.39\%$ & $0.29\%$ \\
Itokawa & $12.8\%$ & $0.9\%$ & $2.59\%$ & $-$      \\
Ida     & $5.6\%$  & $4.3\%$ & $0.11\%$ & $0.05\%$ \\
Gaspra  & $1.8\%$  & $1.5\%$ & $0.17\%$ & $-$      \\
\bottomrule
\end{tabular}
\end{table}

\begin{figure}[tp]
\centering
\includegraphics[width=\textwidth]{fig4.pdf}
\caption{Resolution diagnostics for the two quantities.
(a) Range sensitivity $S$, the peak-to-peak variation over the full mesh range
normalized by the mean. (b) Convergence residual $C$, the fractional change over
the final refinement step. The geopotential dispersion is the more stable diagnostic
on every body by both measures, but the two measures support claims of different
strength: the range advantage is modest on bodies whose slope is already stable,
whereas the convergence residual separates the two cleanly---the dispersion has
reached a plateau on every body, and the slope has not on Ryugu, Eros or
Itokawa.}
\label{fig:diagnostics}
\end{figure}

The convergence residual $C$ draws a sharper and, we argue, more meaningful
distinction. On the three bodies where the slope has not settled---Ryugu, Eros
and Itokawa---the geopotential dispersion has: it changes by $0.60\%$, $0.29\%$ and
nothing at all over the same final refinement step on which the slope still
moves by $8.24\%$, $4.39\%$ and $2.59\%$. On the two bodies where the slope has
settled, the dispersion has settled at least as firmly. There is no body in the
set on which the dispersion is the less converged of the two.

The difference between the two rows of Table~\ref{tab:varsens} is therefore not
merely quantitative. Within the resolutions available, $\potvarn$ appears to
possess a limit that the meshes are approaching, and $\slopeavg$ on the richest
models does not.

\subsection{Why the two diagnostics differ in resolution sensitivity}
\label{sec:whyconverge}

The two diagnostics are computed from the same field, but they depend on the
mesh in fundamentally different ways.

The geopotential of Eq.~\eqref{eq:geopot} is an integral over the body's
volume. Refining the mesh changes the polyhedron's volume and its mass
distribution only slightly, and those changes are what the potential responds
to; a facet subdivided into four contributes very nearly what it contributed
before, because the mass it bounds is nearly unchanged. The dispersion inherits
this behaviour, and both $\potvar$ and $\Ubar$ converge as the polyhedron
converges to the body it approximates.

The slope of Eq.~\eqref{eq:slope} depends on the facet normal $\nhat_{i}$, a
property of the surface at the scale of a single facet. Subdividing a facet
does not leave the normals nearly unchanged: it replaces one direction by
several that differ from it and from each other by whatever relief exists at the
finer scale. There is no reason for that sequence of directions to settle,
because a rough surface continues to present new orientations at every scale at
which it is resolved. This is the familiar situation in which an enclosed volume
converges under refinement while a surface property need not, and it is why the
distinction of Section~\ref{sec:diagnostics} between $S$ and $C$ was worth
drawing: the slope can be made to look stable by choosing a narrow enough range
of meshes, but it does not thereby acquire a limit.

The smoothing experiment of Table~\ref{tab:smooth} supports the same reading
from the other direction. Suppressing the fine relief---removing the very
structure that refinement keeps discovering---lowers the slope substantially,
by $2.9^{\circ}$ on Itokawa and $3.4^{\circ}$ on Eros. It is that structure,
not the body's overall figure, that the slope keeps responding to.

\subsection{What the geopotential dispersion measures}
\label{sec:varmeaning}

Beyond its numerical stability, $\potvarn$ has a direct physical reading. It is
the area-weighted spread of the geopotential over the surface, normalized by
its mean, and therefore measures how far the surface departs from an
equipotential. A body whose surface has fully relaxed into an equipotential
figure would have $\potvarn = 0$; a body far from equilibrium has a large value.

This is what makes it useful beyond bookkeeping. The values in
Table~\ref{tab:varseries} separate the bodies more cleanly than the slope does:
Itokawa at $0.086$ and Gaspra at $0.060$ sit well above Ryugu, Eros and Ida,
which cluster between $0.033$ and $0.036$. Because the quantity is
dimensionless and resolution-stable, such a comparison is meaningful in a way
that a comparison of mean slopes across differently meshed models is not.

The same property underlies its use as an estimator. \citet{kanamaru2019}
minimized a closely related quantity over trial interior density distributions
to infer the internal structure of Itokawa, on the reasoning that a surface
shaped by long-term relaxation should lie close to an equipotential of the true
density field. Any such inversion requires a diagnostic whose variation with
the assumed density is not swamped by its variation with the mesh---which is
exactly the property established here, and which the slope does not possess.
We return to this in Section~\ref{sec:discussion}.

\subsection{Recommendation}
\label{sec:varrecommend}

The results of this section and the last suggest a simple practice. Where a
resolution-independent statement about a body's surface gravity field is
required---for comparison between bodies, between studies, or between shape
models---the normalized geopotential dispersion should be reported, since within
the range tested it converges. The mean slope remains the more directly
interpretable quantity, and nothing here argues against computing it; but it
should be quoted together with the facet count at which it was obtained, and
compared only at matched resolution.

The next section tests both diagnostics under the harder condition that the
shape models themselves differ.
